% appC_L2.tex: two L^2 bounds on a star of arcs.

\section{Two \texorpdfstring{$L^2$}{L2} bounds on a star of arcs}\label{app:L2}

The uniform resolvent bound of Appendix~\ref{app:resolvent} reduces, at the two edges, to two operators on a star of
arcs through a point: a Cauchy operator conjugated by a power weight, and a Cauchy operator with a logarithmic
commutator. Both are dilation invariant, which is what makes the bounds uniform in $n$. We prove them here. The two
lemmas are statements about a fixed geometric configuration; they involve none of the standing assumptions of
Section~\ref{sec:hardedge}, and $|\alpha|<1$ enters only when they are applied with $\beta=\pm\alpha/2$.

\emph{Setting.} A \emph{star} is a set $\Gamma=\bigcup_{j=1}^m\gamma_j([0,R_j])$ of finitely many $C^2$ arcs
$\gamma_j$, parametrised by arc length, with $\gamma_j(0)=0$, pairwise disjoint except at $0$, and with distinct
tangent directions $\gamma_j'(0)=e^{i\theta_j}$. Each arc carries an orientation, and $\Cauchy_\pm$ are the boundary
values of $\Cauchy$ on $\Gamma$ as in \S\ref{sec:tool-cd}. By compactness and transversality there is $C_\Gamma\ge1$ with
\begin{equation}\label{eq:A-geom}
  |\gamma_j'(\rho)-e^{i\theta_j}|\le C_\Gamma\rho,\qquad |\gamma_j(\rho)-\gamma_i(r)|\ge(\rho+r)/C_\Gamma\ (i\ne j),\qquad
  |\gamma_j(\rho)-\gamma_j(r)|\ge|\rho-r|/C_\Gamma .
\end{equation}
A finite union of $C^2$ arcs is a Carleson curve, so $\Cauchy_\pm$ are bounded on $L^2(\Gamma)$ by \S\ref{sec:tool-cd}(a). For
$\lambda>0$ the dilation $(\mathsf D_\lambda\mathfrak h)(z):=\lambda^{-1/2}\mathfrak h(z/\lambda)$ is unitary from $L^2(\Gamma)$ onto
$L^2(\lambda\Gamma)$, and since $\diff s/(s-z)$ is invariant under $(s,z)\mapsto(\lambda s,\lambda z)$, $\mathsf D_\lambda$
intertwines the Cauchy operators of $\Gamma$ and $\lambda\Gamma$. The same holds for the operators below, whose kernels
are $\diff s/(s-z)$ times a function of $(s/|s|,\,z/|z|,\,|z|/|s|)$.

We use one elementary bound repeatedly. Let $\mathfrak k$ be measurable on $(0,\infty)$ with
\[
  \|\mathfrak k\|_*:=\int_0^\infty|\mathfrak k(u)|u^{-1/2}\diff u<\infty,\qquad
  (\mathcal M_{\mathfrak k}\mathfrak h)(r):=\int_0^\infty\mathfrak k(\rho/r)\mathfrak h(\rho)\frac{\diff\rho}r .
\]
Then
\begin{equation}\label{eq:A-minkowski}
  \|\mathcal M_{\mathfrak k}\mathfrak h\|_{L^2(0,\infty)}\le\|\mathfrak k\|_*\,\|\mathfrak h\|_{L^2(0,\infty)} .
\end{equation}
Indeed $(\mathcal M_{\mathfrak k}\mathfrak h)(r)=\int_0^\infty\mathfrak k(u)\mathfrak h(ru)\diff u$, and $\|\mathfrak h(\cdot\,u)\|_{L^2}=u^{-1/2}
\|\mathfrak h\|_{L^2}$, so \eqref{eq:A-minkowski} is Minkowski's integral inequality. Operators
with kernels that are bounded on a bounded set are Hilbert--Schmidt, and we use that without comment.

\begin{lemma}[power-weighted Cauchy operator on a star]\label{lem:A-power}
Let $\Gamma$ be a star and $|\beta|<\frac12$. Then the operators
\[
  \mathfrak h\ \longmapsto\ |z|^{\beta}\,\Cauchy_\pm\bigl(|s|^{-\beta}\mathfrak h\bigr)(z)
\]
are bounded on $L^2(\Gamma)$, with a norm $N_\beta$ that depends only on $\beta$ and $\Gamma$. The same bound holds on
$\lambda\Gamma$ for every $\lambda>0$, on every subset $\Gamma'\subset\lambda\Gamma$ (functions supported on, and
evaluated on, $\Gamma'$), and for the weights $z^{\beta}$, $s^{-\beta}$ (any fixed continuous branches of the argument
on each arc) in place of $|z|^\beta$, $|s|^{-\beta}$, with $N_\beta$ multiplied by $e^{4\pi|\beta|}$.
\end{lemma}

\begin{proof}
We split off the plain Cauchy operator, which is bounded, and compare the remainder, which has a kernel
regular on the diagonal, with its analogue on the star of tangent rays.

Write $T_\pm\mathfrak h:=|z|^\beta\Cauchy_\pm(|s|^{-\beta}\mathfrak h)=\Cauchy_\pm\mathfrak h+R\mathfrak h$, where
\[
  (R\mathfrak h)(z)=\frac1{2\pi i}\int_\Gamma\Bigl(\frac{|z|^\beta}{|s|^\beta}-1\Bigr)\frac{\mathfrak h(s)\diff s}{s-z}.
\]
The kernel of $R$ is bounded near the diagonal of each arc (see below), so $R\mathfrak h$ has no boundary-value
ambiguity, and it suffices to bound $R$. In arc-length coordinates, $z=\gamma_i(r)$, $s=\gamma_j(\rho)$, the kernel of
$R$ is
\[
  \mathfrak r_{ij}(\rho,r)=\frac1{2\pi i}\Bigl(\frac{|\gamma_i(r)|^\beta}{|\gamma_j(\rho)|^\beta}-1\Bigr)
  \frac{\gamma_j'(\rho)}{\gamma_j(\rho)-\gamma_i(r)} .
\]
Its flat analogue, on the star of rays $\{\rho e^{i\theta_j}\}$, is $\mathfrak k_{ij}(\rho/r)/r$ with
\[
  \mathfrak k_{ij}(u):=\frac1{2\pi i}\,\frac{(u^{-\beta}-1)\,e^{i\theta_j}}{ue^{i\theta_j}-e^{i\theta_i}} .
\]
For $i\ne j$ we have $|ue^{i\theta_j}-e^{i\theta_i}|\ge\sin(|\theta_i-\theta_j|/2)(1+u)$, and for $i=j$ the quotient
$(u^{-\beta}-1)/(u-1)$ is continuous at $u=1$. So $|\mathfrak k_{ij}(u)|\le C(1+u^{-\beta})$ for $u\le2$ and
$|\mathfrak k_{ij}(u)|\le C(u^{-\beta}+1)u^{-1}$ for $u\ge2$, and $\|\mathfrak k_{ij}\|_*<\infty$ precisely because
$-\frac12<\beta<\frac12$: near $0$ the integrand is $O(u^{-1/2-\max(\beta,0)})$, near infinity $O(u^{-3/2-\min(\beta,0)})$.
By \eqref{eq:A-minkowski} the flat operator is bounded on $L^2$ of the rays, hence on $L^2$ of the segments $[0,R_j]$.

The difference $\mathfrak r_{ij}(\rho,r)-\mathfrak k_{ij}(\rho/r)/r$ is a sum of two terms. Put
$w_j(\rho):=(|\gamma_j(\rho)|/\rho)^\beta$; by \eqref{eq:A-geom} and $\gamma_j\in C^2$, $w_j$ is $C^1$ on $[0,R_j]$ and
bounded above and below. The first term is
\[
  \frac1{2\pi i}\Bigl(\frac r\rho\Bigr)^\beta\Bigl(\frac{w_i(r)}{w_j(\rho)}-1\Bigr)\frac{\gamma_j'(\rho)}{\gamma_j(\rho)-\gamma_i(r)};
\]
since $|w_i(r)-w_j(\rho)|\le C|\rho-r|$ for $i=j$ and $\le C(\rho+r)$ for $i\ne j$, \eqref{eq:A-geom} bounds it by
$C(r/\rho)^\beta$. The second term is $\frac1{2\pi i}((r/\rho)^\beta-1)\,\mathfrak e_{ij}(\rho,r)$ with
\[
  \mathfrak e_{ij}(\rho,r):=\frac{\gamma_j'(\rho)}{\gamma_j(\rho)-\gamma_i(r)}-\frac{e^{i\theta_j}}{\rho e^{i\theta_j}-re^{i\theta_i}} .
\]
We claim $|\mathfrak e_{ij}|\le C$. For $i\ne j$ both denominators are at least $(\rho+r)/C$, and the numerator of the
difference, $\gamma_j'(\rho)(\rho e^{i\theta_j}-re^{i\theta_i})-e^{i\theta_j}(\gamma_j(\rho)-\gamma_i(r))$, is
$O(\rho(\rho+r))+O(\rho^2+r^2)$ by \eqref{eq:A-geom}. For $i=j$ the flat term is $1/(\rho-r)$, and
$\gamma_j(\rho)-\gamma_j(r)=(\rho-r)(\gamma_j'(\rho)+O(\rho-r))$, so $\gamma_j'(\rho)/(\gamma_j(\rho)-\gamma_j(r))=
(\rho-r)^{-1}(1+O(\rho-r))$, which gives $\mathfrak e_{jj}=O(1)$. Here the factor $\gamma_j'(\rho)$, i.e.\ the arc-length
element $\diff s=\gamma_j'(\rho)\diff\rho$, must be kept: without it the difference is only $O(r/|\rho-r|)$ on a single
arc. Hence the difference kernel is at most $C(1+(r/\rho)^\beta)$, and
$\iint_{[0,R]^2}(1+(r/\rho)^{2\beta})\diff\rho\diff r<\infty$ for $|\beta|<\frac12$: it is Hilbert--Schmidt.

So $R$ is bounded on $L^2(\Gamma)$ (arc length makes the parametrisation unitary), and so is $T_\pm$. Dilation invariance
gives the bound on $\lambda\Gamma$. Restriction to $\Gamma'$ is a compression. Finally $z^\beta=|z|^\beta e^{i\beta\arg z}$
with $|\arg z|\le2\pi$, so the complex weights differ from the absolute ones by multiplication operators of norm at most
$e^{2\pi|\beta|}$ on either side.
\end{proof}

\begin{lemma}[logarithmic commutator on a star]\label{lem:A-log}
Let $\Gamma$ be a star and fix on each arc a continuous branch $\mathrm{l}(z)=\log|z|+i\arg z$ of the logarithm. Then the
operator
\[
  (L\mathfrak h)(z):=\frac1{2\pi i}\int_\Gamma\bigl(\mathrm l(z)-\mathrm l(s)\bigr)\frac{\mathfrak h(s)\diff s}{s-z}
\]
is bounded on $L^2(\Gamma)$, with the same norm on $\lambda\Gamma$ for every $\lambda>0$ and on subsets.
\end{lemma}

\begin{proof}
The kernel is regular on the diagonal, since $\mathrm l$ is Lipschitz in arc length away from $0$ and the factor
$\mathrm l(z)-\mathrm l(s)$ vanishes at $s=z$; we split it into the part with $\log|z|-\log|s|$ and the part with
$\arg z-\arg s$.

\emph{The modulus part.} On the star of rays its kernel is $\mathfrak k_{ij}(\rho/r)/r$ with
\[
  \mathfrak k_{ij}(u)=-\frac1{2\pi i}\,\frac{\log u\;e^{i\theta_j}}{ue^{i\theta_j}-e^{i\theta_i}} .
\]
This has no singularity at
$u=1$ even for $i=j$ ($\log u/(u-1)\to1$), and $|\mathfrak k_{ij}(u)|u^{-1/2}$ is $O(|\log u|u^{-1/2})$ at $0$ and
$O(|\log u|u^{-3/2})$ at infinity, so $\|\mathfrak k_{ij}\|_*<\infty$ and \eqref{eq:A-minkowski} applies. On the curved star
the difference from the flat kernel is
\[
  \frac1{2\pi i}\Bigl[(\lambda_i(r)-\lambda_j(\rho))\frac{\gamma_j'(\rho)}{\gamma_j(\rho)-\gamma_i(r)}
  +\log\frac r\rho\;\mathfrak e_{ij}(\rho,r)\Bigr],
\]
with $\lambda_j(\rho):=\log(|\gamma_j(\rho)|/\rho)$, a $C^1$ function on
$[0,R_j]$. As in the proof of Lemma~\ref{lem:A-power} this is $O(1+|\log(r/\rho)|)$, and
$\iint_{[0,R]^2}(1+\log^2(r/\rho))\diff\rho\diff r<\infty$.

\emph{The argument part.} Its kernel is $\frac{i}{2\pi i}(\arg\gamma_i(r)-\arg\gamma_j(\rho))\gamma_j'(\rho)/(\gamma_j(\rho)-
\gamma_i(r))$. For $i=j$, $\arg\gamma_j(\rho)=\theta_j+O(\rho)$ is $C^1$ on $[0,R_j]$, so by \eqref{eq:A-geom} the kernel is
bounded. For $i\ne j$ the numerator is bounded and the denominator is at least $(\rho+r)/C_\Gamma$, so the kernel is at most
$C/(\rho+r)=C\,\mathfrak k(\rho/r)/r$ with $\mathfrak k(u)=(1+u)^{-1}$, $\|\mathfrak k\|_*=\pi$.

Both parts are bounded on $L^2(\Gamma)$. The kernel depends on $(s,z)$ only through $s/|s|$, $z/|z|$ and $|z|/|s|$,
so the operator commutes with $\mathsf D_\lambda$ and the bound holds on $\lambda\Gamma$.
\end{proof}

\begin{remark}[the condition $|\beta|<\frac12$]\label{rem:mellin}
For $\Gamma$ a union of rays the curvature terms vanish, and Lemma~\ref{lem:A-power} reduces to the boundedness of the
Mellin multiplier $\int_0^\infty u^{w-1}\diff u/(u-e^{i\phi_{ij}})$ on the line $\re w=\frac12-\beta$; on a single
ray its principal value is $-\pi\cot(\pi w)$ (by \cite[(5.12.3), (5.5.3)]{DLMF}, see \S\ref{sec:tool-dlmf}), which is
bounded on that line exactly when $0<\frac12-\beta<1$. This explains the role of the condition $|\beta|<\frac12$.
\end{remark}
